% Hafta 10 -- alıştırma gövdesi (ders notu ve ayrı alıştırma kağıdı tarafından \input edilir)
\section{Çözümlü Alıştırmalar}
Her soruyu önce kendiniz deneyin, sonra çözümü okuyun.

\begin{alistirma}
\(\sqrt[3]{-8}+\sqrt[4]{16}-\sqrt[5]{-32}\) işleminin sonucu kaçtır?
\end{alistirma}
\begin{cozum}
Tek dereceli köklerde negatif içerik serbesttir:
\(\sqrt[3]{-8}=-2\) \ (\((-2)^3=-8\)), \quad \(\sqrt[4]{16}=2\) \ (\(2^4=16\)), \quad \(\sqrt[5]{-32}=-2\) \ (\((-2)^5=-32\)).
\[-2+2-(-2)=-2+2+2=2\]
\end{cozum}

\begin{alistirma}
\(\sqrt{50}+\sqrt8-\sqrt{32}\) işleminin sonucu nedir?
\end{alistirma}
\begin{cozum}
Önce her kökü sadeleştirelim (en büyük tam kare çarpanı ayıralım):\\
\(\sqrt{50}=\sqrt{25\cdot 2}=5\sqrt2\), \quad \(\sqrt8=\sqrt{4\cdot 2}=2\sqrt2\), \quad \(\sqrt{32}=\sqrt{16\cdot 2}=4\sqrt2\).
\[5\sqrt2+2\sqrt2-4\sqrt2=(5+2-4)\sqrt2=3\sqrt2\]
\end{cozum}

\begin{alistirma}
\(\dfrac{\sqrt{48}}{\sqrt3}+\dfrac{\sqrt{45}}{\sqrt5}\) işleminin sonucu kaçtır?
\end{alistirma}
\begin{cozum}
Aynı dereceli kökler bölünürken içler bölünür:
\[\frac{\sqrt{48}}{\sqrt3}=\sqrt{\frac{48}{3}}=\sqrt{16}=4,\qquad \frac{\sqrt{45}}{\sqrt5}=\sqrt{\frac{45}{5}}=\sqrt9=3,\qquad 4+3=7\]
\end{cozum}

\begin{alistirma}
\(\dfrac{4}{\sqrt5-1}\) kesrinin paydasını rasyonel yapınız.
\end{alistirma}
\begin{cozum}
Paydanın eşleniği \(\sqrt5+1\). Pay ve payda bununla çarpılır; payda iki kare farkıdır:
\((\sqrt5-1)(\sqrt5+1)=5-1=4\).
\[\frac{4(\sqrt5+1)}{4}=\sqrt5+1\]
\end{cozum}

\begin{alistirma}
\(\sqrt{2\sqrt[3]{4}}\) ifadesini tek kök olarak yazınız.
\end{alistirma}
\begin{cozum}
Önce \(2\)'yi küp köke alalım: \(2\sqrt[3]{4}=\sqrt[3]{2^3\cdot 4}=\sqrt[3]{32}\).\\
Kökün kökü: \(\sqrt{\sqrt[3]{32}}=\sqrt[2\cdot 3]{32}=\sqrt[6]{32}\). \quad (Üslü yazımla: \(2\cdot 2^{\frac23}=2^{\frac53}\), \ \(\left(2^{\frac53}\right)^{\frac12}=2^{\frac56}=\sqrt[6]{32}\).)
\end{cozum}

\begin{alistirma}
\(a=\sqrt[3]{3}\), \ \(b=\sqrt2\), \ \(c=\sqrt[6]{10}\) sayılarını küçükten büyüğe sıralayınız.
\end{alistirma}
\begin{cozum}
Kök dereceleri \(3\), \(2\), \(6\); EKOK \(=6\). Dereceleri eşitleyelim (içerinin üssü de aynı sayıyla çarpılır):
\[a=\sqrt[6]{3^2}=\sqrt[6]{9},\qquad b=\sqrt[6]{2^3}=\sqrt[6]{8},\qquad c=\sqrt[6]{10}\]
\(8<9<10\) \(\Rightarrow\) \(b<a<c\).
\end{cozum}

\begin{alistirma}
\(\sqrt{8+2\sqrt{15}}-\sqrt{8-2\sqrt{15}}\) işleminin sonucu nedir?
\end{alistirma}
\begin{cozum}
Toplamı \(8\), çarpımı \(15\) olan sayılar \(5\) ve \(3\):
\[\sqrt{8+2\sqrt{15}}=\sqrt5+\sqrt3,\qquad \sqrt{8-2\sqrt{15}}=\sqrt5-\sqrt3\quad(\text{büyük olan önce})\]
\[(\sqrt5+\sqrt3)-(\sqrt5-\sqrt3)=2\sqrt3\]
\end{cozum}

\begin{alistirma}
\(\sqrt{3x+4}=5\) \ ve \ \(\sqrt[3]{2^{y+1}}=4\) olduğuna göre \(x+y\) kaçtır?
\end{alistirma}
\begin{cozum}
1. denklem: iki tarafın karesi alınır: \(3x+4=25\Rightarrow 3x=21\Rightarrow x=7\). \ (Kontrol: \(\sqrt{25}=5\).)\\
2. denklem: kesirli üs ile \(2^{\frac{y+1}{3}}=2^2\Rightarrow\dfrac{y+1}{3}=2\Rightarrow y=5\). \ (Kontrol: \(\sqrt[3]{2^6}=\sqrt[3]{64}=4\).)
\[x+y=7+5=12\]
\end{cozum}

\kendinizoncesi
\section{Alıştırmalar (Kendiniz Çözün)}

\begin{enumerate}[label=\textbf{\arabic*.},leftmargin=*,itemsep=6pt]
\item \(\sqrt{0{,}04}+\sqrt{1{,}44}-\sqrt{0{,}25}\) işleminin sonucu kaçtır?
\item \(\sqrt[4]{81}+\sqrt[3]{-64}\) işleminin sonucu kaçtır?
\item \(\sqrt{18}\), \(\sqrt{75}\) ve \(\sqrt{128}\) sayılarını \(a\sqrt b\) biçiminde yazınız.
\item \(3\sqrt2+5\sqrt2-7\sqrt2\) işleminin sonucu nedir?
\item \(\sqrt3\cdot\sqrt6\cdot\sqrt2\) çarpımı kaçtır?
\item \(9^{\frac32}-32^{\frac25}\) işleminin sonucu kaçtır?
\item \(\sqrt{(2-\sqrt7)^2}\) ifadesinin eşiti nedir?
\item \(\sqrt{60}\) hangi iki ardışık tam sayı arasındadır?
\item \(\dfrac{6}{\sqrt3}\) kesrinin paydasını rasyonel yapınız.
\item \(\sqrt[3]{3}\cdot\sqrt3\) çarpımını tek kök olarak yazınız.
\item \(a=2\sqrt5\), \(b=3\sqrt2\), \(c=\sqrt{19}\) sayılarını küçükten büyüğe sıralayınız.
\item \(\dfrac{1}{\sqrt6-\sqrt5}\) kesrinin paydasını rasyonel yapınız.
\item \(\dfrac{2}{\sqrt3+1}+\dfrac{2}{\sqrt3-1}\) işleminin sonucu nedir?
\item \(\sqrt{6-2\sqrt5}\) ifadesinin eşiti nedir?
\item \(x=\sqrt2+1\) olduğuna göre \(x^2-2x\) kaçtır? (Kareyi dağılma özelliğiyle açınız.)
\item \(x=\sqrt{30+\sqrt{30+\sqrt{30+\cdots}}}\) olduğuna göre \(x\) kaçtır?
\end{enumerate}

\needspace{14\baselineskip}
\subsection*{Cevap Anahtarı}
{\small
\begin{multicols}{2}
\begin{enumerate}[label=\textbf{\arabic*.},leftmargin=*,itemsep=1pt]
\item \(0{,}2+1{,}2-0{,}5=0{,}9\).
\item \(3+(-4)=-1\).
\item \(3\sqrt2\); \ \(5\sqrt3\); \ \(\sqrt{64\cdot 2}=8\sqrt2\).
\item \((3+5-7)\sqrt2=\sqrt2\).
\item \(\sqrt{36}=6\).
\item \((\sqrt9)^3-(\sqrt[5]{32})^2=27-4=23\).
\item \(\sqrt7>2\Rightarrow|2-\sqrt7|=\sqrt7-2\).
\item \(49<60<64\Rightarrow 7\) ile \(8\) arası.
\item \(\dfrac{6\sqrt3}{3}=2\sqrt3\).
\item \(\sqrt[6]{3^2}\cdot\sqrt[6]{3^3}=\sqrt[6]{243}\).
\item \(\sqrt{20},\ \sqrt{18},\ \sqrt{19}\Rightarrow b<c<a\).
\item Payda \(6-5=1\Rightarrow\sqrt6+\sqrt5\).
\item \((\sqrt3-1)+(\sqrt3+1)=2\sqrt3\).
\item \(5+1=6\), \(5\cdot 1=5\Rightarrow\sqrt5-1\).
\item \(x^2=2+\sqrt2+\sqrt2+1=3+2\sqrt2\); \ \(3+2\sqrt2-2\sqrt2-2=1\).
\item \(x(x-1)=30=6\cdot 5\Rightarrow x=6\).
\end{enumerate}
\end{multicols}}
